% !TeX root = ../../金融经济学_王江_讲义及习题.tex
% 知识点：references/金融经济学/金融经济学-王江_261013.pdf，印刷页 179–192（PDF 页 194–207）
% 习题：references/金融经济学/史上最全最完整的金融经济学答案(王江版).pdf，页 62–67（印刷页 61–66）

\lecture{在均值-方差偏好下的投资组合选择}

\ifdefined\WJHandoutKnowledge

上一章在一般偏好和不完全证券结构下，均衡选择依赖完整的支付分布。本章增加条件，使参与者只按收益分布的\textbf{均值}与\textbf{方差}选择组合；这把组合问题化为协方差矩阵上的二次优化。

\begin{tcolorbox}[knowledge,title=12.1 均值--方差偏好]
记终值财富为 $\widetilde w$，均值 $\bar w=\mathrm E[\widetilde w]$，标准差 $\sigma_w=\sqrt{\operatorname{Var}(\widetilde w)}$。若参与者的排序可由
\[
 \mathrm E[u(\widetilde w)]=v(\bar w,\sigma_w)
 \tag{12.1}
\]
完全表示，则称具有\textbf{均值--方差偏好}。这需要对效用函数或可选收益分布施加条件，而非一般期望效用的自动性质。

\textbf{定理 12.1}\quad 若 $u(w)=w-\tfrac12 a w^2$，则
\[
 \mathrm E[u(\widetilde w)]
 =\bar w-\tfrac12a(\bar w^2+\sigma_w^2).
\]
因此二次效用具有均值--方差表示。但 $u'(w)=1-aw$ 只在 $w<1/a$ 时为正；若财富越过该区间，二次效用会违反不满足性，使用时须限制财富域。

\textbf{定理 12.2}\quad 若可交易证券的联合支付为正态分布，则每个线性组合仍为正态分布，其分布完全由均值和方差决定。因此，只要期望效用存在，参与者对这些组合的排序具有均值--方差表示；不要求效用本身二次型。

\textbf{定理 12.3}\quad 对联合正态支付及递增、严格凹的效用，$v$ 对均值严格递增、对标准差严格递减；$v$ 对均值凹、对标准差凹，且具有相应的联合凹性。证明把财富写成 $\bar w+\sigma_w Z$，$Z\sim N(0,1)$：
\[
 v_1=\mathrm E[u'(\bar w+\sigma_wZ)]>0,
 \quad v_2=\mathrm E[u'(\bar w+\sigma_wZ)Z]<0.
\]
后一个符号可由 $Z$ 对称性与 $u'$ 递减得到。二阶导数的凹性由 $u''<0$ 及 Cauchy--Schwarz 不等式得到。正态分布有负财富尾部；若效用只定义在正财富域，必须另作定义域处理，不能直接套用。
\end{tcolorbox}

\begin{tcolorbox}[knowledge,title=12.2 无风险证券缺席时的前沿]
设 $N$ 只风险证券的净收益率向量为 $\widetilde r$，均值 $\bar r$，正定协方差矩阵 $\Sigma$。权重 $z$ 满足 $\mathbf1^{\mathsf T}z=1$，组合均值和方差为
\[
 \bar r_p=z^{\mathsf T}\bar r,\qquad
 \sigma_p^2=z^{\mathsf T}\Sigma z.
\]
在给定目标收益 $\bar r_p$ 下求最小方差，得到\textbf{均值--方差前沿}（MVF）组合：
\[
 \min_z\ \tfrac12 z^{\mathsf T}\Sigma z,
 \quad\text{s.t. }z^{\mathsf T}\bar r=\bar r_p,
 \quad z^{\mathsf T}\mathbf1=1.
 \tag{12.2}
\]
令 $a=\bar r^{\mathsf T}\Sigma^{-1}\bar r$、$b=\bar r^{\mathsf T}\Sigma^{-1}\mathbf1$、$c=\mathbf1^{\mathsf T}\Sigma^{-1}\mathbf1$、$d=ac-b^2>0$，一阶条件解为
\[
 z_p=\frac1d\Sigma^{-1}\bigl[(c\bar r_p-b)\bar r+(a-b\bar r_p)\mathbf1\bigr].
 \tag{12.3}
\]
$d>0$ 要求各风险证券的预期收益不全相同；若 $d=0$，目标收益约束退化，应单独处理。正定的 $\Sigma$ 保证解唯一。

把解写成 $z_p=z_0+\bar r_p(z_1-z_0)$，可见\textbf{定理 12.5}：任意 MVF 组合可由 $z_0,z_1$ 两个 MVF 组合形成；\textbf{推论 12.1--12.2}：任意两个不同的 MVF 组合足以生成整条前沿，它们的线性组合仍在前沿。这就是无风险证券缺席时的\textbf{两基金分离}（定理 12.6--12.7）。组合权重允许卖空；若另有权重非负约束，前沿会分段，公式须改写。

前沿在 $(\sigma,\bar r)$ 平面是一条双曲线的边界。顶点是\textbf{最小方差组合}（MVP），其均值为 $b/c$，权重为 $\Sigma^{-1}\mathbf1/c$。前沿的上半支在均值--方差偏好下有效；下半支被同方差、更高均值的组合支配。

\textbf{定理 12.8}\quad 对任意组合 $p$，$\operatorname{Cov}(\widetilde r_p,\widetilde r_{\rm MVP})=\sigma_{\rm MVP}^2$。证明：MVP 的一阶条件给出 $\Sigma z_{\rm MVP}=\mathbf1/c$，故 $z_p^{\mathsf T}\Sigma z_{\rm MVP}=1/c$。\textbf{定理 12.9}进一步指出，对非 MVP 的前沿组合 $p$，存在一个均值--方差前沿上的零协方差组合 $zcp$，用作后面的定价基准。

\textbf{定理 12.10}\quad 若 $p$ 是非 MVP 的前沿组合，则任意组合 $q$ 满足
\[
 \bar r_q-\bar r_{zcp}
 =\beta_{q,p}(\bar r_p-\bar r_{zcp}),\qquad
 \beta_{q,p}=\frac{\operatorname{Cov}(\widetilde r_q,\widetilde r_p)}{\sigma_p^2}.
 \tag{12.6}
\]
证明将 $\Sigma z_p=\lambda_r\bar r+\lambda_1\mathbf1$ 两边左乘任意权重 $z_q^{\mathsf T}$，再分别取 $q=p,zcp$ 消去乘子。这里是\textbf{零贝塔定价式}，其基准收益不一定等于无风险利率。
\end{tcolorbox}

\begin{tcolorbox}[knowledge,title=12.4 存在无风险证券时]
令无风险净利率为 $r_F$，风险证券的预期超额收益向量 $\eta=\bar r-r_F\mathbf1$。风险证券权重 $z$ 不必合计为 1，余下资金投向无风险证券。若 $\eta^{\mathsf T}\Sigma^{-1}\eta>0$，给定组合预期超额收益 $\bar\eta_p$，最小方差解是
\[
 z_p=\frac{\bar\eta_p}{\eta^{\mathsf T}\Sigma^{-1}\eta}\Sigma^{-1}\eta.
 \tag{12.7}
\]
将 $\bar\eta_p$ 取为使风险证券权重和为 1 的值，得到\textbf{切点组合}
\[
 z_T=\frac{\Sigma^{-1}\eta}{\mathbf1^{\mathsf T}\Sigma^{-1}\eta},
\]
前提是分母非零。所有此类前沿组合是无风险证券与 $T$ 的线性组合（\textbf{定理 12.11}）；这就是一基金加无风险证券的分离。若 $r_F$ 位于风险证券前沿的适当位置，切点 $T$ 在有效前沿上（\textbf{定理 12.12}）。

无风险证券与 $T$ 在 $(\sigma,\bar r)$ 平面形成\textbf{资本市场线}。风险组合 $q$ 的 Sharpe 比为 $(\bar r_q-r_F)/\sigma_q$；在允许自由借贷和卖空的上述条件下，切点组合具有最高 Sharpe 比。

\textbf{定理 12.13}\quad 若 $p$ 是有无风险证券时的均值--方差有效组合，任意组合 $q$ 满足
\[
 \bar r_q-r_F=\beta_{q,p}(\bar r_p-r_F),\qquad
 \beta_{q,p}=\frac{\operatorname{Cov}(\widetilde r_q,\widetilde r_p)}{\sigma_p^2}.
 \tag{12.9}
\]
特别地，取 $p=T$ 得切点组合的贝塔定价式（12.10）。这里的 $p$ 是\emph{有效组合}；下一章才在均衡条件下识别它与\emph{市场组合}的关系。
\end{tcolorbox}

\begin{tcolorbox}[summary,title=本章小结]
二次效用或联合正态支付等条件可使组合排序只依赖均值与方差。无风险证券缺席时，前沿由两个风险基金生成，并有零贝塔定价关系；存在无风险证券时，前沿由切点风险组合与无风险证券生成，风险溢价可由相对于有效组合的贝塔表示。上述分析以协方差矩阵正定、相关组合可交易及适当的定义域为条件。
\end{tcolorbox}

\fi

% !TeX root = ../../金融经济学_王江_讲义及习题.tex
% 12_均值方差偏好 课后习题（由章节文件拆分）

\begin{problem}
\problemno{12.1} 令交易证券的价格向量为 $p$、总收益率向量为 $\tilde{x}=\begin{bmatrix}\tilde{x}_{1} \\ \dots \\ \tilde{x}_{N}\end{bmatrix}$。假设 $\tilde{x}=\iota+\tilde{r}$ 服从均值为 $\iota+\bar{r}$、协方差阵为 $\Sigma$ 的正态分布。对于初始财富为 $w$ 的参与者，他的 1 期财富是 $\tilde{w}=wz^{\top}\tilde{x}$，其中 $z$ 是他的权重向量（$\iota^{\top}z=1$）。假设参与者的绝对风险厌恶系数恒为 $a$。
\begin{enumerate}[label=(\alph*)]
    \item 计算他的间接效用函数
    \[
        v=\mathrm{E}[u(\tilde{w})]
    \]
    其中 $u$ 是他的直接效用函数。
    \item 证明间接效用函数具有均值--方差表示，即它可以表示成 $\bar{w}=\mathrm{E}[\tilde{w}]$ 和 $\sigma_{w}^{2}=\mathrm{Var}[\tilde{w}]$ 的函数，即 $v(\bar{w},\sigma_{w})$。
    \item 证明均值方差效用函数 $v(\bar{w},\sigma_{w})$ 满足如下条件：$\partial_{1}v>0$，$\partial_{2}v<0$ 且 $\partial_{11}^{2}v<0$，$\partial_{22}^{2}v<0$ 且 $v$ 是凹的。
\end{enumerate}
\end{problem}
\begin{solution}
不妨假设效用函数为 $u(w)=-e^{-aw}$，那么 $v=\mathrm{E}[u(\tilde{w})]=\mathrm{E}[-e^{-a\tilde{w}}]$。因为 $\tilde{w}$ 是正态分布的，因而
\[
    v=-e^{-a\mathrm{E}\tilde{w}+\frac12 a^{2}\mathrm{Var}\tilde{w}}=-e^{-a\bar{w}+\frac12 a^{2}\sigma_{w}^{2}}=v(\bar{w},\sigma_{w})
\]
即间接效用函数具有均值方差表示，对 $v$ 求导可得
\[
    \begin{aligned}
        \partial_{1}v&=e^{-a\bar{w}+\frac12 a^{2}\sigma_{w}^{2}}a>0\\
        \partial_{2}v&=-e^{-a\bar{w}+\frac12 a^{2}\sigma_{w}^{2}}a^{2}\sigma_{w}<0\\
        \partial_{11}v&=e^{-a\bar{w}+\frac12 a^{2}\sigma_{w}^{2}}(-a^{2})<0\\
        \partial_{22}v&=-e^{-a\bar{w}+\frac12 a^{2}\sigma_{w}^{2}}(a^{4}\sigma_{w}^{2}+a^{2})<0\\
        \partial_{12}v&=e^{-a\bar{w}+\frac12 a^{2}\sigma_{w}^{2}}a^{3}\sigma_{w},\qquad
        \partial_{11}v\,\partial_{22}v-(\partial_{12}v)^{2}=v^{2}(\bar{w},\sigma_{w})a^{2}>0
    \end{aligned}
\]
因而，$v$ 还是凹的。
\end{solution}

\begin{problem}
\problemno{12.2} 考虑证券 1 和 2，它们的收益率分别是 $\tilde{r}_{1}$ 和 $\tilde{r}_{2}$。它们的波动率分别是 $\sigma_{1}$ 和 $\sigma_{2}$，并且它们的相关系数是 $\rho$。
\begin{enumerate}[label=(\alph*)]
    \item 假设 $\rho\ne 1$。用波动率度量风险，证明分散化即用证券形成组合可以降低风险。（也就是说，总是存在波动率小于两只单个证券的波动率的组合。）
    \item 找出收益率具有最小波动率的组合。
    \item 证明当 $\rho=1$ 时，分散化不能降低风险。
    \item 假设 $\sigma_{1}=\sigma_{2}$。证明无论 $\rho$ 的值是多少，等权重组合的风险是最小的。
\end{enumerate}
\end{problem}
\begin{solution}
\begin{enumerate}[label=(\alph*)]
    \item 设组合权重为 $z=\begin{bmatrix}z_{1} \\ 1-z_{1}\end{bmatrix}$，那么 $\tilde{r}_{p}=z_{1}\tilde{r}_{1}+(1-z_{1})\tilde{r}_{2}$，
    \[
        \sigma_{p}^{2}=z_{1}^{2}\sigma_{1}^{2}+2z_{1}(1-z_{1})\rho\sigma_{1}\sigma_{2}+(1-z_{1})^{2}\sigma_{2}^{2}
    \]
    对 $z_{1}$ 求导，$\partial\sigma_{p}^{2}/\partial z_{1}=0$ 给出
    \[
        z_{1}\sigma_{1}^{2}+(1-2z_{1})\rho\sigma_{1}\sigma_{2}-(1-z_{1})\sigma_{2}^{2}=0
    \]
    因而全局最小方差权重为
    \[
        z_{1}^{*}=\frac{\sigma_{2}^{2}-\rho\sigma_{1}\sigma_{2}}{\sigma_{1}^{2}-2\rho\sigma_{1}\sigma_{2}+\sigma_{2}^{2}}
    \]
    代回并配方，
    \[
        \sigma_{p}^{2}=\sigma_{2}^{2}-\frac{(\sigma_{2}^{2}-\rho\sigma_{1}\sigma_{2})^{2}}{\sigma_{1}^{2}-2\rho\sigma_{1}\sigma_{2}+\sigma_{2}^{2}}=\sigma_{1}^{2}-\frac{(\sigma_{1}^{2}-\rho\sigma_{1}\sigma_{2})^{2}}{\sigma_{1}^{2}-2\rho\sigma_{1}\sigma_{2}+\sigma_{2}^{2}}
    \]
    当 $\rho\ne 1$ 时分子为正，故 $\sigma_{p}<\min(\sigma_{1},\sigma_{2})$，分散化可以降低风险；
    \item 即上一问中的 $z^{*}$；
    \item 当 $\rho=1$ 时，$\sigma_{p}^{2}=(\sigma_{1}z_{1}+\sigma_{2}(1-z_{1}))^{2}\in[\min(\sigma_{1}^{2},\sigma_{2}^{2}),\,\max(\sigma_{1}^{2},\sigma_{2}^{2})]$，$z_{1}\in[0,1]$，即分散化不能降低风险。
    \item 当 $\sigma_{1}=\sigma_{2}$ 时，$z_{1}=(\sigma_{2}^{2}-\rho\sigma_{1}\sigma_{2})/(\sigma_{1}^{2}-2\rho\sigma_{1}\sigma_{2}+\sigma_{2}^{2})=\tfrac12$，即无论 $\rho$ 多少，等权重组合的风险最小。
\end{enumerate}
\end{solution}

\begin{problem}
\problemno{12.3*} 假设市场上有 $N$ 只证券，它们收益率波动率均为 $\sigma$，而协方差为 $c=\mathrm{Cov}[\tilde{r}_{i},\tilde{r}_{j}]$，$\forall\, i,j=1,\dots,N$ 且 $i\ne j$。
\begin{enumerate}[label=(\alph*)]
    \item 证明由 $N$ 只证券组成的等权组合的收益率具有最小波动率。
    \item 证明当 $N\to\infty$ 时等权重组合的方差趋近于 $c$。这和单个证券的波动率 $\sigma$ 无关。
    \item 解释上面得到的结论。
\end{enumerate}
\end{problem}
\begin{solution}
\begin{enumerate}[label=(\alph*)]
    \item 假设组合权重为 $z=\begin{bmatrix}z_{1} \\ z_{2} \\ \dots \\ z_{N}\end{bmatrix}$，那么组合的方差为
    \[
        \sigma_{p}^{2}=\mathrm{Var}(z_{1}\tilde{r}_{1}+z_{2}\tilde{r}_{2}+\dots+z_{N}\tilde{r}_{N})=\sigma^{2}\sum_{i}z_{i}^{2}+c\sum_{i\ne j}z_{i}z_{j}
    \]
    利用 $\sum_{i}z_{i}=1$ 改写为 $\sigma_{p}^{2}=(\sigma^{2}-c)\sum_{i}z_{i}^{2}+c$。拉格朗日函数为
    \[
        \mathcal{L}=(\sigma^{2}-c)\sum_{i}z_{i}^{2}+\lambda\Bigl(1-\sum_{i}z_{i}\Bigr)
    \]
    一阶条件 $2(\sigma^{2}-c)z_{i}=\lambda$（$\sigma^{2}-c>0$）迫使 $z_{i}$ 全相等，因此
    \[
        z_{i}=\frac{1}{N},\quad\forall i
    \]
    \item 把 $z_{i}=\dfrac{1}{N}$ 代入组合方差可得 $\sigma_{p}^{2}=\dfrac{1}{N}\sigma^{2}+\bigl(1-\dfrac{1}{N}\bigr)c\to c$，$N\to\infty$；
    \item 当 $N\to\infty$ 时，证券间不相关的风险都被分散掉了，只剩下了其中相关的部分，因而最终的组合风险只与协方差项有关。
\end{enumerate}
\end{solution}

\begin{problem}
\problemno{12.4} 考虑风险证券 1 和 2，它们的收益率为 $\tilde{r}_{1}$ 和 $\tilde{r}_{2}$。它们的期望收益率和协方差阵分别为 $\bar{r}_{1}$，$\bar{r}_{2}$，以及 $\Sigma=\begin{bmatrix}\sigma_{1}^{2} & \sigma_{12} \\ \sigma_{12} & \sigma_{2}^{2}\end{bmatrix}$。
\begin{enumerate}[label=(\alph*)]
    \item 推导由这两只证券构成的均值--方差有效前沿。
    \item 证明这两只证券都在前沿上。
    \item 考虑如下三种特别形式的前沿：$\rho_{12}=-1,\,0,\,1$，其中 $\rho_{12}$ 是两只证券收益率之间的相关系数。在 $\bar{r}$--$\sigma$ 平面画出这三种情形下的组合前沿。
\end{enumerate}
\end{problem}
\begin{solution}
\begin{enumerate}[label=(\alph*)]
    \item 为简便起见，假设 $\bar{r}_{1}\ne\bar{r}_{2}$；给定组合的期望收益率 $\bar{r}_{p}$ 有 $z\bar{r}_{1}+(1-z)\bar{r}_{2}=\bar{r}_{p}$，因而 $z=\dfrac{\bar{r}_{p}-\bar{r}_{2}}{\bar{r}_{1}-\bar{r}_{2}}$，此时，组合的波动率为
    \[
        \begin{aligned}
            \sigma_{p}^{2}&=\frac{1}{(\bar{r}_{1}-\bar{r}_{2})^{2}}\bigl[(\sigma_{1}^{2}+\sigma_{2}^{2}-2\sigma_{12})\bar{r}_{p}^{2}-(2\sigma_{1}^{2}\bar{r}_{2}+2\sigma_{2}^{2}\bar{r}_{1}-2(\bar{r}_{1}+\bar{r}_{2})\sigma_{12})\bar{r}_{p}\\
            &\qquad+(\sigma_{1}^{2}\bar{r}_{2}^{2}+\sigma_{2}^{2}\bar{r}_{1}^{2}-2\bar{r}_{1}\bar{r}_{2}\sigma_{12})\bigr]
        \end{aligned}
    \]
    $(\bar{r}_{p},\sigma_{p})$ 即为均值--方差前沿组合；
    \item 当 $\bar{r}_{p}=\bar{r}_{1}$ 时，前沿组合即为证券 1；当 $\bar{r}_{p}=\bar{r}_{2}$ 时，前沿组合即为证券 2；
    \item 当 $\rho_{12}=-1$ 时
    \[
        \sigma_{p}=\left|\frac{(\sigma_{1}+\sigma_{2})\bar{r}_{p}-(\sigma_{1}\bar{r}_{2}+\sigma_{2}\bar{r}_{1})}{\bar{r}_{1}-\bar{r}_{2}}\right|
    \]
    当 $\rho_{12}=1$ 时
    \[
        \sigma_{p}=\left|\frac{(\sigma_{1}-\sigma_{2})\bar{r}_{p}-(\sigma_{1}\bar{r}_{2}+\sigma_{2}\bar{r}_{1})}{\bar{r}_{1}-\bar{r}_{2}}\right|
    \]
    在这两种情形下，前沿组合由两条起点一直的射线构成。当 $\rho_{12}=0$ 时，
    \[
        \sigma_{p}^{2}=\frac{1}{(\bar{r}_{1}-\bar{r}_{2})^{2}}\bigl[(\sigma_{1}^{2}+\sigma_{2}^{2}-2\sigma_{12})\bar{r}_{p}^{2}-(2\sigma_{1}^{2}\bar{r}_{2}+2\sigma_{2}^{2}\bar{r}_{1})\bar{r}_{p}+(\sigma_{1}^{2}\bar{r}_{2}^{2}+\sigma_{2}^{2}\bar{r}_{1}^{2})\bigr]
    \]
    此时的前沿组合是 $\bar{r}$--$\sigma$ 平面中的一条抛物线。
\end{enumerate}
\end{solution}

\begin{problem}
\problemno{12.5} 令市场上风险证券的收益率向量为 $\tilde{r}$，它的均值为 $\bar{r}$，它的协方差阵为 $\Sigma$，$\Sigma$ 是正定的，而 $\iota$ 为同阶的元素全为 1 的向量。另外，令 $a=\iota^{\top}\Sigma^{-1}\bar{r}$，$b=\bar{r}^{\top}\Sigma^{-1}\bar{r}$，$c=\iota^{\top}\Sigma^{-1}\iota$ 且 $d=bc-a^{2}$。证明 $d=0$ 等价于 $\bar{r}=(b/a)\iota$，即所有风险证券具有相同的期望收益 $(b/a)$。
\end{problem}
\begin{solution}
当 $d=0$ 时，$bd=(a\bar{r}-b\iota)^{\top}\Sigma^{-1}(a\bar{r}-b\iota)=0$，因为 $\Sigma$，$\Sigma^{-1}$ 都是正定的，因而 $a\bar{r}-b\iota=0$，$\bar{r}=\dfrac{b}{a}\iota$，即所有风险证券具有相同的期望收益 $\dfrac{b}{a}$。
\end{solution}

\begin{problem}
\problemno{12.6} 证明最小方差组合（MVP）的期望收益率为 $a/c$，这里的 $a$ 和 $c$ 由练习 12.5 给定。
\end{problem}
\begin{solution}
我们知道最小方差组合的期望收益率为 $r_{\mathrm{MVP}}=-\dfrac{(z_{1}-z_{0})^{\top}\Sigma^{-1}z_{0}}{(z_{1}-z_{0})^{\top}\Sigma^{-1}(z_{1}-z_{0})}$，$z_{0}$，$z_{1}$ 由正文 2.3 节定义，把 $z_{0}$，$z_{1}$ 代入得
\[
    \begin{aligned}
        r_{\mathrm{MVP}}
        &=-\frac{d^{-2}(c\Sigma^{-1}\bar{r}-a\Sigma^{-1}\iota)^{\top}\Sigma(b\Sigma^{-1}\iota-a\Sigma^{-1}\bar{r})}{d^{-2}(c\Sigma^{-1}\bar{r}-a\Sigma^{-1}\iota)^{\top}\Sigma(c\Sigma^{-1}\bar{r}-a\Sigma^{-1}\iota)}\\
        &=-\frac{a(a^{2}-bc)}{c(bc-a^{2})}=\frac{a}{c}
    \end{aligned}
\]
\end{solution}

\begin{problem}
\problemno{12.7} 考虑 $\bar{r}=(b/a)\iota$ 的特殊情形。导出均值--方差前沿。在 $\bar{r}$--$\sigma$ 平面上画出这个组合前沿。
\end{problem}
\begin{solution}
当期望收益率都相等时，所有组合的期望收益率都是常数，因而均值--方差前沿组合退化成为一点即最小方差组合，此时的方差为 $\sigma_{\mathrm{MVP}}^{2}=\dfrac{1}{\iota^{\top}\Sigma^{-1}\iota}$，对应的组合权重为 $z_{\mathrm{MVP}}=\dfrac{\Sigma^{-1}\iota}{\iota^{\top}\Sigma^{-1}\iota}$。
\end{solution}

\begin{problem}
\problemno{12.8**} 令证券收益率向量为 $\tilde{r}$，它的均值为 $\bar{r}$、协方差阵为 $\Sigma$。在 $\bar{r}$--$\sigma$ 平面上，任一组合 $p$ 由一点 $(\bar{r}_{p},\sigma_{p})$ 给定。令所有可能组合的集合为 $A=\{(\bar{r}_{p},\sigma_{p}),\forall\, p\}$。
\begin{enumerate}[label=(\alph*)]
    \item 证明 $A$ 是闭集。
    \item $A$ 一定是凸集吗？若否，给出反例。
\end{enumerate}
\end{problem}
\begin{solution}
\begin{enumerate}[label=(\alph*)]
    \item 假设有一序列 $\bigl\{(r_{p}^{(n)},\sigma_{p}^{(n)})\bigr\}_{1}^{\infty}$，其极限为 $(r_{p},\sigma_{p})$。那么存在 $\theta^{(n)}\in S^{N-1}=\{\theta\in\mathrm{R}^{N}:\theta^{\top}\iota=1\}$ 使得
    \[
        \theta^{(n)\top}\bar{r}=r_{p}^{(n)},\qquad\theta^{(n)\top}\Sigma\theta^{(n)}=\sigma_{p}^{(n)2}
    \]
    显然，$\theta^{(n)}$ 是 $S^{N-1}$ 中的有界序列，存在收敛子列 $\{\theta^{(k_{n})}\}_{1}^{\infty}$，记其极限为 $\theta\in S^{N-1}$，那么
    \[
        \begin{aligned}
            \theta^{\top}\bar{r}&=\lim_{n}\theta^{(k_{n})\top}\bar{r}=\lim_{n}r_{p}^{(k_{n})}=r_{p}\\
            \theta^{\top}\Sigma\theta&=\lim_{n}\theta^{(k_{n})\top}\Sigma\theta^{(k_{n})}=\lim_{n}\sigma_{p}^{(k_{n})2}=\sigma_{p}^{2}
        \end{aligned}
    \]
    因而 $(r_{p},\sigma_{p})\in A$，$A$ 是闭集；
    \item 不一定。当 $N=2$ 且两只证券的期望收益不同、收益率不完全相关时，卖空下的组合在 $\bar{r}$--$\sigma$ 平面上只扫出均值--方差前沿，即一条双曲线。双曲线不是凸集：取前沿上期望收益不同的两点，连接它们的线段会落到前沿左侧，对应的 $(\bar{r},\sigma)$ 不可由任何组合达到。
\end{enumerate}
\end{solution}

\begin{problem}
\problemno{12.9**} 令 $p$ 为 MVF 组合，$zcp$ 是与 $p$ 的协方差为零的 MVF。
\begin{enumerate}[label=(\alph*)]
    \item 证明对任意证券 $i$，我们有
    \[
        \tilde{r}_{i}=\alpha+\beta_{i,zcp}\tilde{r}_{zcp}+\beta_{i,p}\tilde{r}_{p}+\tilde{u}_{i}
    \]
    其中 $\mathrm{Cov}[\tilde{r}_{zcp},\tilde{u}_{i}]=0$，$\mathrm{Cov}[\tilde{r}_{p},\tilde{u}_{i}]=0$，$\mathrm{E}[\tilde{u}_{i}]=0$，$\beta_{i,zcp}=\mathrm{Cov}[\tilde{r}_{i},\tilde{r}_{zcp}]/\mathrm{Var}[\tilde{r}_{zcp}]$ 且 $\beta_{i,p}=\mathrm{Cov}[\tilde{r}_{i},\tilde{r}_{p}]/\mathrm{Var}[\tilde{r}_{p}]$。
    \item 证明 $\beta_{i,zcp}=1-\beta_{i,p}$。
    \item 证明 $\alpha=0$ 及 $\bar{r}_{i}=(1-\beta_{i,p})\bar{r}_{zcp}+\beta_{i,p}\bar{r}_{p}$。
\end{enumerate}
\end{problem}
\begin{solution}
\begin{enumerate}[label=(\alph*)]
    \item 令 $\tilde{r}_{i}=\alpha+\beta_{i,1}\tilde{r}_{zcp}+\beta_{i,2}\tilde{r}_{p}+\tilde{u}_{i}$，由
    \[
        \mathrm{Cov}(\tilde{u}_{i},\tilde{r}_{p})=\mathrm{Cov}(\tilde{r}_{i}-\alpha-\beta_{i,1}\tilde{r}_{zcp}-\beta_{i,2}\tilde{r}_{p},\tilde{r}_{p})=0
    \]
    可得 $\beta_{i,2}=\beta_{i,p}$；同样地有 $\beta_{i,1}=\beta_{i,zcp}$，由 $\mathrm{E}\tilde{u}_{i}=0$ 可以得到 $\alpha=\bar{r}_{i}-\beta_{i,p}\bar{r}_{p}-\beta_{i,zcp}\bar{r}_{zcp}$；
    \item 由文中 (12.6) 式有
    \[
        \begin{aligned}
            \bar{r}_{i}-\bar{r}_{zcp}&=\beta_{i,p}(\bar{r}_{p}-\bar{r}_{zcp})\\
            \bar{r}_{i}-\bar{r}_{p}&=\beta_{i,zcp}(\bar{r}_{zcp}-\bar{r}_{p})
        \end{aligned}
    \]
    两式相减可以得到 $\beta_{i,p}+\beta_{i,zcp}=1$；
    \item 另外，由 (12.6) 式有
    \[
        \bar{r}_{i}=\bar{r}_{zcp}+\beta_{i,p}(\bar{r}_{p}-\bar{r}_{zcp})=\beta_{i,p}\bar{r}_{p}+(1-\beta_{i,p})\bar{r}_{zcp}
    \]
    因此 $\alpha=0$。
\end{enumerate}
\end{solution}

\begin{problem}
\problemno{12.10} 证明 MVE 组合的凸组合（即权重非负的组合）仍然是 MVE 组合。
\end{problem}
\begin{solution}
首先，我们知道 MVE 组合的组合是 MVF 组合；另外，由于是凸组合，因而所得 MVF 组合在均值--方差前沿的上半部，即为 MVE 组合。
\end{solution}

\begin{problem}
\problemno{12.11} 当存在一只收益率为 $r_{F}$ 的无风险债券时，记 $\bar{\eta}=\bar{r}-r_{F}\iota$ 是风险证券的超额收益向量。证明 $\bar{\eta}^{\top}\Sigma^{-1}\bar{\eta}=0$ 等价于 $\bar{r}=r_{F}\iota$。
\end{problem}
\begin{solution}
因为 $\Sigma^{-1}$ 是正定的，
\[
    \bar{\eta}^{\top}\Sigma^{-1}\bar{\eta}=0\Leftrightarrow\bar{\eta}=\bar{r}-r_{F}\iota=0\Leftrightarrow\bar{r}=r_{F}\iota
\]
\end{solution}

\begin{problem}
\problemno{12.12} 当存在无风险证券时，切点组合的期望超额收益为正等价于 $\bar{r}_{\mathrm{mvp}}-r_{F}>0$，即最小方差组合的期望收益为正。
\end{problem}
\begin{solution}
当切点组合的期望收益率为正时，
\[
    \begin{aligned}
        \iota^{\top}\Sigma^{-1}\bar{\eta}>0
        &\Leftrightarrow\iota^{\top}\Sigma^{-1}(\bar{r}-r_{F}\iota)=\iota^{\top}\Sigma^{-1}\bar{r}-r_{F}\iota^{\top}\Sigma^{-1}\iota=\iota^{\top}\Sigma^{-1}\iota\Bigl(\frac{\iota^{\top}\Sigma^{-1}\bar{r}}{\iota^{\top}\Sigma^{-1}\iota}-r_{F}\Bigr)\\
        &=\iota^{\top}\Sigma^{-1}\iota\Bigl(\frac{a}{c}-r_{F}\Bigr)=\iota^{\top}\Sigma^{-1}\iota(\bar{r}_{\mathrm{mvp}}-r_{F})>0\\
        &\Leftrightarrow\bar{r}_{\mathrm{mvp}}-r_{F}>0.
    \end{aligned}
\]
\end{solution}

\begin{problem}
\problemno{12.13*} 当存在无风险证券时，证明切点组合的 Sharpe 比最大。
\end{problem}
\begin{solution}
考虑任意组合 Sharpe 比的平方
\[
    \frac{z^{\top}\bar{r}+(1-z^{\top}\iota)r_{F}-r_{F}}{z^{\top}\Sigma z}=\frac{z^{\top}\bar{\eta}}{z^{\top}\Sigma z}
\]
其中 $z$ 是风险组合的权重，最大化的一阶条件是
\[
    \bar{\eta}\,z^{\top}\Sigma z=z\,\bar{\eta}^{\top}\Sigma z
\]
可以验证 $z_{T}=\dfrac{1}{\iota^{\top}\Sigma^{-1}\bar{\eta}}\Sigma^{-1}\bar{\eta}$ 即为其解，因而切点组合的 Sharpe 比最大。
\end{solution}


